\documentclass[12pt]{article}
\usepackage[utf8]{inputenc}
\usepackage[T1]{fontenc}
\usepackage[margin=1in]{geometry}
\usepackage{setspace}
\usepackage{xcolor}
\usepackage[colorlinks=true, urlcolor=blue!60!black, linkcolor=blue!60!black]{hyperref}
\setlength{\parindent}{0em}
\setlength{\parskip}{0.6em}
\renewcommand{\baselinestretch}{1.0}
\begin{document}
\begin{center}
{\LARGE \textit{United States v. Comstock}, 560 U.S. 126 (2010)\par}
\vspace{0.5em}
{\large Case Brief --- Two-Page Edition\par}
\end{center}

\section{Memory Jogger}
By 7--2, the Court held that the Necessary and Proper Clause lets Congress civilly commit sexually dangerous federal prisoners beyond their release dates, as a narrow incident of running federal prisons and acting as a responsible custodian (\href{https://supreme.justia.com/cases/federal/us/560/126/}{560 U.S. 126}).

\section{Detailed Case Facts}
Under 18 U.S.C. \S~4248, enacted in 2006, a federal court may commit a ``sexually dangerous'' person in federal custody if the government proves by clear and convincing evidence past sexually violent conduct or child molestation, a serious mental illness, and resulting serious difficulty refraining from such conduct. The Attorney General must try to transfer custody to the person's State.

\section{Procedural History}
The government sought commitment of Graydon Comstock and four others. The District Court and the Fourth Circuit held that Congress lacked Article I power; the Supreme Court granted certiorari on that question.

\section{Judicial Votes}
Reversed, 7--2. Breyer wrote for Roberts, Stevens, Ginsburg, and Sotomayor; Kennedy and Alito concurred in the judgment; Thomas, joined in most part by Scalia, dissented.

\section{Holding}
Section 4248 is ``a means that is rationally related to the implementation of a constitutionally enumerated power'' and falls within the Necessary and Proper Clause.

\section{Analysis of Opinions}
Breyer relied on five considerations: the Clause's breadth under \textit{McCulloch}; the long history of federal commitment of mentally ill prisoners; the government's role as a responsible custodian; the statute's accommodation of state interests; and its narrow reach to people already in federal custody. Kennedy warned against borrowing ``rational basis'' rhetoric, and Alito demanded a ``substantial link,'' found here because federal imprisonment severs inmates' state ties. Thomas argued that the power behind a conviction is exhausted when the sentence ends.

\section{Examples: Future Applications}
\textbf{Same-Side:} holding a contagious inmate in isolation past release until his State takes custody is valid. \textbf{Opposite-Side:} committing past state felons never in federal custody, via a firearm's interstate travel, is invalid. \textbf{Fence-Sitter:} new registration duties for federal offenders released decades ago.

\section{Critique}
Critics echo Thomas: stacking implied powers on implied powers leaves no principled stopping point. Kennedy, Alito, and Roberts in \textit{NFIB} recast ``proper'' as an independent federalism check.

\section{Key Quotations}
\begin{quote}
``Taken together, these considerations lead us to conclude that the statute is a `necessary and proper' means of exercising the federal authority that permits Congress to create federal criminal laws, to punish their violation, to imprison violators, to provide appropriately for those imprisoned, and to maintain the security of those who are not imprisoned but who may be affected by the federal imprisonment of others.''
\end{quote}
\noindent --- Justice Breyer, Opinion of the Court, 560 U.S. at 149.

\end{document}
